\documentclass{article}
\usepackage[paperwidth=60cm,paperheight=120cm,margin=0pt]{geometry}
\usepackage{fontspec}\setmainfont{Latin Modern Roman}
\usepackage{xeCJK}\setCJKmainfont{Microsoft YaHei}
\usepackage{amsmath, amssymb, siunitx}
\usepackage[american]{circuitikz}
\ctikzset{bipoles/length=1.2cm}
% 不用 standalone：它自家的裁边在 xelatex 下不生效，长图会被 letter 页面截掉尾部。
% 改用 preview 把页面尺寸严格设成内容尺寸（自动裁边）。
\usepackage[active,tightpage]{preview}
\PreviewEnvironment{circuitikz}
\setlength\PreviewBorder{8pt}
\pagestyle{empty}
\begin{document}

\begin{circuitikz}
\small
\node[anchor=east, font=\bfseries\small] at (4.55,0.00) {MCU\ /\ 电源};
\node[anchor=west, font=\bfseries\small] at (4.95,0.00) {网络名};
\node[anchor=west, font=\bfseries\small] at (12.45,0.00) {连接器针脚};
\node[anchor=west, font=\bfseries\small] at (12.30,-0.60) {TB6612FNG 电机驱动模块};
\node[anchor=west, font=\scriptsize\itshape] at (12.30,-0.98) {方向线全部走 PB12~PB15（不需要关 JTAG）；VM 走大电流焊盘};
\node[anchor=east, font=\small] at (4.55,-1.84) {PB6\ \textperiodcentered\ TIM4\_CH1};
\draw (4.85,-1.84) -- (12.30,-1.84);
\node[anchor=south, font=\scriptsize\ttfamily] at (8.57,-1.79) {PWM\_L};
\node[anchor=west, font=\small] at (12.48,-1.84) {1\,--\,PWMA};
\node[anchor=east, font=\small] at (4.55,-2.58) {PB7\ \textperiodcentered\ TIM4\_CH2};
\draw (4.85,-2.58) -- (12.30,-2.58);
\node[anchor=south, font=\scriptsize\ttfamily] at (8.57,-2.53) {PWM\_R};
\node[anchor=west, font=\small] at (12.48,-2.58) {2\,--\,PWMB};
\node[anchor=east, font=\small] at (4.55,-3.32) {PB12};
\draw (4.85,-3.32) -- (12.30,-3.32);
\node[anchor=south, font=\scriptsize\ttfamily] at (8.57,-3.27) {AIN1\_L};
\node[anchor=west, font=\small] at (12.48,-3.32) {3\,--\,AIN1};
\node[anchor=east, font=\small] at (4.55,-4.06) {PB13};
\draw (4.85,-4.06) -- (12.30,-4.06);
\node[anchor=south, font=\scriptsize\ttfamily] at (8.57,-4.01) {AIN2\_L};
\node[anchor=west, font=\small] at (12.48,-4.06) {4\,--\,AIN2};
\node[anchor=east, font=\small] at (4.55,-4.80) {PB14};
\draw (4.85,-4.80) -- (12.30,-4.80);
\node[anchor=south, font=\scriptsize\ttfamily] at (8.57,-4.75) {BIN1\_R};
\node[anchor=west, font=\small] at (12.48,-4.80) {5\,--\,BIN1};
\node[anchor=east, font=\small] at (4.55,-5.54) {PB15};
\draw (4.85,-5.54) -- (12.30,-5.54);
\node[anchor=south, font=\scriptsize\ttfamily] at (8.57,-5.49) {BIN2\_R};
\node[anchor=west, font=\small] at (12.48,-5.54) {6\,--\,BIN2};
\node[anchor=east, font=\small] at (4.55,-6.28) {PA11\ \textperiodcentered\ 硬件 10k 上拉 + 跳线可选常使能};
\draw (4.85,-6.28) -- (12.30,-6.28);
\node[anchor=south, font=\scriptsize\ttfamily] at (8.57,-6.23) {STBY};
\node[anchor=west, font=\small] at (12.48,-6.28) {7\,--\,STBY};
\node[anchor=east, font=\small] at (4.55,-7.02) {5V\ \textperiodcentered\ 模块自带 3.3V LDO 供 3V3};
\draw (4.85,-7.02) -- (12.30,-7.02);
\node[anchor=south, font=\scriptsize\ttfamily] at (8.57,-6.97) {5V};
\node[anchor=west, font=\small] at (12.48,-7.02) {8\,--\,VCC};
\node[anchor=east, font=\small] at (4.55,-7.76) {GND};
\draw (4.85,-7.76) -- (12.30,-7.76);
\node[anchor=south, font=\scriptsize\ttfamily] at (8.57,-7.71) {GND};
\node[anchor=west, font=\small] at (12.48,-7.76) {9\,--\,GND};
\node[anchor=east, font=\small] at (4.55,-8.50) {12V\ \textperiodcentered\ 大电流焊盘 + 470µF};
\draw (4.85,-8.50) -- (12.30,-8.50);
\node[anchor=south, font=\scriptsize\ttfamily] at (8.57,-8.45) {VM};
\node[anchor=west, font=\small] at (12.48,-8.50) {10\,--\,VM};
\draw[rounded corners=2pt] (12.30,-1.42) rectangle (18.40,-8.80);
\node[anchor=west, font=\bfseries\small] at (12.30,-9.40) {左电机 JST-XH 6P 座（现成插头）};
\node[anchor=west, font=\scriptsize\itshape] at (12.30,-9.78) {同一个座里同时走电机电流与编码器信号 —— 线序必须实测复核；A/B 接反只需软件换向};
\node[anchor=east, font=\small] at (4.55,-10.64) {左电机\ \textperiodcentered\ 大电流};
\draw (4.85,-10.64) -- (12.30,-10.64);
\node[anchor=south, font=\scriptsize\ttfamily] at (8.57,-10.59) {MOT\_L+};
\node[anchor=west, font=\small] at (12.48,-10.64) {1\,--\,M+};
\node[anchor=east, font=\small] at (4.55,-11.38) {左电机\ \textperiodcentered\ 大电流};
\draw (4.85,-11.38) -- (12.30,-11.38);
\node[anchor=south, font=\scriptsize\ttfamily] at (8.57,-11.33) {MOT\_L-};
\node[anchor=west, font=\small] at (12.48,-11.38) {2\,--\,M−};
\node[anchor=east, font=\small] at (4.55,-12.12) {3V3\ \textperiodcentered\ 跳线可选 5V};
\draw (4.85,-12.12) -- (12.30,-12.12);
\node[anchor=south, font=\scriptsize\ttfamily] at (8.57,-12.07) {ENC\_VCC};
\node[anchor=west, font=\small] at (12.48,-12.12) {3\,--\,编码器 VCC};
\node[anchor=east, font=\small] at (4.55,-12.86) {GND};
\draw (4.85,-12.86) -- (12.30,-12.86);
\node[anchor=south, font=\scriptsize\ttfamily] at (8.57,-12.81) {GND};
\node[anchor=west, font=\small] at (12.48,-12.86) {4\,--\,GND};
\node[anchor=east, font=\small] at (4.55,-13.60) {PA0\ \textperiodcentered\ TIM2\_CH1};
\draw (4.85,-13.60) -- (12.30,-13.60);
\node[anchor=south, font=\scriptsize\ttfamily] at (8.57,-13.55) {ENC\_L\_A};
\node[anchor=west, font=\small] at (12.48,-13.60) {5\,--\,A};
\node[anchor=east, font=\small] at (4.55,-14.34) {PA1\ \textperiodcentered\ TIM2\_CH2};
\draw (4.85,-14.34) -- (12.30,-14.34);
\node[anchor=south, font=\scriptsize\ttfamily] at (8.57,-14.29) {ENC\_L\_B};
\node[anchor=west, font=\small] at (12.48,-14.34) {6\,--\,B};
\draw[rounded corners=2pt] (12.30,-10.22) rectangle (18.40,-14.64);
\node[anchor=west, font=\bfseries\small] at (12.30,-15.24) {右电机 JST-XH 6P 座（现成插头）};
\node[anchor=east, font=\small] at (4.55,-16.08) {右电机\ \textperiodcentered\ 大电流};
\draw (4.85,-16.08) -- (12.30,-16.08);
\node[anchor=south, font=\scriptsize\ttfamily] at (8.57,-16.03) {MOT\_R+};
\node[anchor=west, font=\small] at (12.48,-16.08) {1\,--\,M+};
\node[anchor=east, font=\small] at (4.55,-16.82) {右电机\ \textperiodcentered\ 大电流};
\draw (4.85,-16.82) -- (12.30,-16.82);
\node[anchor=south, font=\scriptsize\ttfamily] at (8.57,-16.77) {MOT\_R-};
\node[anchor=west, font=\small] at (12.48,-16.82) {2\,--\,M−};
\node[anchor=east, font=\small] at (4.55,-17.56) {3V3\ \textperiodcentered\ 跳线可选 5V};
\draw (4.85,-17.56) -- (12.30,-17.56);
\node[anchor=south, font=\scriptsize\ttfamily] at (8.57,-17.51) {ENC\_VCC};
\node[anchor=west, font=\small] at (12.48,-17.56) {3\,--\,编码器 VCC};
\node[anchor=east, font=\small] at (4.55,-18.30) {GND};
\draw (4.85,-18.30) -- (12.30,-18.30);
\node[anchor=south, font=\scriptsize\ttfamily] at (8.57,-18.25) {GND};
\node[anchor=west, font=\small] at (12.48,-18.30) {4\,--\,GND};
\node[anchor=east, font=\small] at (4.55,-19.04) {PA6\ \textperiodcentered\ TIM3\_CH1};
\draw (4.85,-19.04) -- (12.30,-19.04);
\node[anchor=south, font=\scriptsize\ttfamily] at (8.57,-18.99) {ENC\_R\_A};
\node[anchor=west, font=\small] at (12.48,-19.04) {5\,--\,A};
\node[anchor=east, font=\small] at (4.55,-19.78) {PA7\ \textperiodcentered\ TIM3\_CH2};
\draw (4.85,-19.78) -- (12.30,-19.78);
\node[anchor=south, font=\scriptsize\ttfamily] at (8.57,-19.73) {ENC\_R\_B};
\node[anchor=west, font=\small] at (12.48,-19.78) {6\,--\,B};
\draw[rounded corners=2pt] (12.30,-15.66) rectangle (18.40,-20.08);
\node[anchor=west, font=\scriptsize\itshape] at (0,-20.98) {网络名 = 嘉立创 EDA 里的网络标号：两边同名即同一根线。电源（5V / 3V3 / GND / 12V）同样用网络标号表示，不画共用母线。逐针定义、电气约束与改动规则见 电气连接规范.md；引脚契约见 docs/引脚分配表.md。};
\end{circuitikz}
\end{document}
